% ============================================================================
%  AI Investment, EU vs US --- Project Documentation
%  Compile with:  pdflatex documentation.tex   (run twice for references)
%  Font: Palatino (mathpazo, with matching math).
% ============================================================================
\documentclass[11pt,a4paper]{article}

% ---- Fonts: Palatino for text and math -------------------------------------
\usepackage[T1]{fontenc}
\usepackage[utf8]{inputenc}
\usepackage{mathpazo}            % Palatino text + matching maths (Hermann Zapf)
\linespread{1.05}                % Palatino reads better with a touch more leading
\usepackage{microtype}

% ---- Layout ----------------------------------------------------------------
\usepackage[a4paper,margin=2.6cm]{geometry}
\usepackage{booktabs}
\usepackage{array}
\usepackage{amsmath,amssymb}
\usepackage{graphicx}
\graphicspath{{figures/}{./}}
\usepackage{caption}
\captionsetup{font=small,labelfont=bf}
\usepackage{enumitem}
\setlist{itemsep=2pt,topsep=3pt}
\usepackage{xcolor}
\definecolor{codebg}{gray}{0.96}
\definecolor{rule}{gray}{0.5}

% ---- Code / paths listing style --------------------------------------------
\usepackage{listings}
\lstset{
  basicstyle=\ttfamily\small,
  backgroundcolor=\color{codebg},
  frame=single, rulecolor=\color{rule},
  breaklines=true, columns=fullflexible,
  showstringspaces=false, keepspaces=true,
  xleftmargin=4pt, xrightmargin=4pt, framesep=4pt,
}

\usepackage[hidelinks,colorlinks=true,linkcolor=black!70!blue,
            urlcolor=black!60!blue,citecolor=black]{hyperref}

% ---- Convenience macros ----------------------------------------------------
\newcommand{\code}[1]{\texttt{#1}}
\newcommand{\dataset}[1]{\textsc{#1}}
% Robust figure inclusion: shows a placeholder if the PNG is absent.
\newcommand{\projfig}[3]{% {width-fraction}{file-without-ext}{caption}
  \begin{figure}[htbp]\centering
    \IfFileExists{figures/#2.png}
      {\includegraphics[width=#1\linewidth]{#2}}
      {\fbox{\parbox{0.78\linewidth}{\centering\itshape figure not found ---
        run the pipeline (\code{02\_make\_figures.py}) to generate it.}}}
    \caption{#3}\label{fig:#2}
  \end{figure}}

\title{\textbf{AI Investment, EU vs US}\\[2pt]
       \large Data, Pipeline, and the ``Adopt-vs-Fund'' Gap}
\author{Project documentation}
\date{\today}

\begin{document}
\maketitle
\thispagestyle{empty}

\begin{abstract}
\noindent
This document explains, step by step, the empirical project behind a single
claim: \emph{Europe is becoming a consumer of artificial intelligence rather
than a producer of it.} Two directly observable facts support it. First, European
firms increasingly \emph{use} AI---adoption has risen steeply since 2023. Second,
almost none of the private capital that \emph{builds} frontier AI is European:
investment is overwhelmingly American, and the gap is widening. The document
covers (i) the reproducible two-stage software pipeline, (ii) every data series
used, with an explicit justification of its source, frequency, coverage, units,
and treatment, and (iii) the descriptive charts that carry the argument. The
project is \emph{descriptive by design}: the thesis is an accounting comparison
of adoption against investment, not a causal mechanism, so it runs no
econometrics. Reported figures come from a clean end-to-end run; the adoption
series is pulled live and may be revised by the provider, and the investment
figures should be refreshed from the latest report before citation.
\end{abstract}

\tableofcontents
\bigskip
\hrule
\bigskip

% ============================================================================
\section{Overview and research question}\label{sec:overview}
% ============================================================================

\paragraph{Thesis.} A common framing of the EU's position in AI fixates on the
venture-capital gap with the United States. This project sharpens the point into
a testable contrast: Europe \emph{adopts} AI without \emph{funding} it. The claim
breaks into three legs, each mapped to a chart:
\begin{enumerate}[label=(\Alph*)]
  \item \textbf{Investment.} The money building AI is concentrated in the United
        States to an extraordinary degree, and the lead is growing, not shrinking
        (the headline leg).
  \item \textbf{Adoption.} Yet European \emph{use} of AI is rising fast---the EU
        is a large and growing customer for AI it largely did not build.
  \item \textbf{Within-Europe divides.} Both the using and (what little) funding
        happen unevenly: a north-west/south-east geographic split, and a
        large-firm/SME split.
\end{enumerate}

\paragraph{Why descriptive, not econometric.} The sibling precautionary-saving
project devotes most of its length to isolating a \emph{latent} causal channel
(does uncertainty \emph{drive} saving?), because there co-movement is not
causation. Here the situation is different and simpler: both halves of the thesis
are \emph{directly measured} quantities---a share of firms using AI, and a flow
of dollars invested---so the argument is an accounting comparison, not an
identification problem. There is no hidden variable to recover and nothing to
instrument. The honest altitude is therefore a small set of clean, well-sourced
charts; adding a regression would dress up a descriptive claim as a causal one.
What the project owes the reader instead is rigour about \emph{data provenance
and comparability}, which Section~\ref{sec:data} supplies.

\paragraph{What the documentation covers.} Section~\ref{sec:pipeline} describes
the reproducible pipeline. Section~\ref{sec:data} justifies every data choice and
states the two comparability caveats that matter. Section~\ref{sec:descriptive}
walks through the five charts. Sections~\ref{sec:interpretation}--%
\ref{sec:limits} synthesise and qualify the findings. The appendix lists the file
tree, exact commands, and a data dictionary.

% ============================================================================
\section{The reproducible pipeline}\label{sec:pipeline}
% ============================================================================

The project is split into two scripts by \emph{function}, so that each stage can
be run, audited, and re-run independently. Data collection (slow, network-bound,
non-deterministic across vintages) is separated from figure rendering (fast,
deterministic, offline). There is deliberately no third stage: per
Section~\ref{sec:overview} the analysis is descriptive.

\begin{center}
\renewcommand{\arraystretch}{1.25}
\begin{tabular}{@{}p{4.4cm} p{4.4cm} p{5.8cm}@{}}
\toprule
\textbf{Stage / script} & \textbf{Reads} & \textbf{Writes} \\
\midrule
\code{01\_collect\_data.py} & Eurostat (live API) + AI Index figures &
  \code{data/*.csv} (tidy series) \\
\code{02\_make\_figures.py} & \code{data/*.csv} &
  \code{figures/*.png} (charts A--B3) \\
\bottomrule
\end{tabular}
\end{center}

\paragraph{One-command reproduction.} The wrapper \code{run\_all.sh} installs the
pinned dependencies and runs the two stages in order:
\begin{lstlisting}
bash run_all.sh        # pip install, then run 01 -> 02 (classic Python)
\end{lstlisting}

\paragraph{Classic Python, no virtual environment.} The project runs directly on
a standard system Python ($\ge 3.11$); there is nothing to create or activate.
\code{run\_all.sh} installs the requirements into the chosen interpreter
(\code{python3} by default; override with the \code{PYTHON} variable) and then
runs the two stages. Dependencies are pinned in \code{requirements.txt} so the
charts are reproducible across machines. If the default \code{python3} is a
Homebrew/``externally-managed'' build that refuses a global \code{pip install},
point \code{PYTHON} at a python.org interpreter or add \code{--user} to the pip
line.

\paragraph{Network robustness.} The adoption data is pulled through the keyless
\code{eurostat} package. Each pull is wrapped in \code{try/except}: a failed
source prints a clear message and leaves any CSV from a previous run untouched,
so a transient outage never corrupts the rest of the dataset. The latest survey
year is detected automatically from the returned table, so a new Eurostat release
flows through to the country bar and the firm-size split with no code change.

% ============================================================================
\section{Data}\label{sec:data}
% ============================================================================

This section answers the ``why this data'' question. For each series we state the
\emph{source} and exact dataset code, the \emph{frequency} and why it is
appropriate, the \emph{coverage}, the \emph{units}, and any \emph{transformation}
and \emph{caveat}. The adoption series is free and keyless; the investment
figures are public but, as explained below, are not exposed through any free API.

% ----------------------------------------------------------------------------
\subsection{Adoption: AI use by EU enterprises}
% ----------------------------------------------------------------------------
\begin{description}[leftmargin=1.4em,style=nextline]
\item[Source / code] Eurostat, \dataset{isoc\_eb\_ai} (``Artificial intelligence
  by size class of enterprise''), indicator \code{E\_AI\_TANY} (enterprises using
  at least one of the surveyed AI technologies), unit \code{PC\_ENT}
  (\% of enterprises), activity \code{C10-S951\_X\_K} (all activities excluding
  the financial sector). Pulled live.
\item[Frequency / coverage] Annual, but the module is not run every year:
  available years are 2021, 2023, 2024, 2025 (no 2022). The latest year drives
  the country and firm-size charts; all years feed the trend chart.
\item[Units] Percentage of enterprises with 10 or more persons employed
  (size class \code{GE10}) using $\ge 1$ AI technology. This is an
  \emph{extensive} measure---whether a firm uses any AI at all---not an intensity
  or spending measure.
\item[Why this exact indicator] \dataset{isoc\_eb\_ai} contains dozens of
  indicators (by purpose, by technology type, by barrier, by provenance).
  \code{E\_AI\_TANY} is the single headline ``does the firm use any AI'' switch
  and the one Eurostat itself reports as the adoption rate; the collector filters
  to it explicitly rather than blending indicators.
\item[Caveat: series break] \emph{Eurostat broadened the list of surveyed AI
  technologies between the 2023 and 2024 waves.} The EU-27 rate jumps from
  $8.1\%$ (2023) to $13.5\%$ (2024); part of that step is the wider definition,
  not pure real growth. The $2024\!\to\!2025$ step ($13.5\%\!\to\!19.9\%$) is on
  the consistent definition and is clean. The trend chart shades the break and
  the caption flags it.
\end{description}

% ----------------------------------------------------------------------------
\subsection{Adoption by firm size (leg C)}
% ----------------------------------------------------------------------------
\begin{description}[leftmargin=1.4em,style=nextline]
\item[Source] The same \dataset{isoc\_eb\_ai}/\code{E\_AI\_TANY} series,
  geography \code{EU27\_2020}, broken out by the size classes \code{10-49}
  (small), \code{50-249} (medium), and \code{GE250} (large).
\item[Why] AI adoption is sharply increasing in firm size; the small/large
  contrast is the cleanest within-EU statement of the ``who actually uses it''
  question, and is robust (a single country-free EU aggregate).
\end{description}

% ----------------------------------------------------------------------------
\subsection{Investment: private AI funding by region (legs A)}
% ----------------------------------------------------------------------------
\begin{description}[leftmargin=1.4em,style=nextline]
\item[Source] Stanford HAI \emph{AI Index} reports (2024, 2025, 2026 editions);
  underlying data assembled by \dataset{Quid}. Current-USD private investment.
\item[Coverage / units] Three views are used: (i) the single-year cross-country
  bar for the latest year, 2025; (ii) the US/Europe/China anchor-year path
  2023--2025; (iii) the cumulative 2013--2025 sum. All in billions of current US
  dollars of \emph{private} investment.
\item[Why hardcoded, with provenance] The underlying series is proprietary
  (Quid) and is \emph{not} exposed through any free API; the public form is the
  figures printed in the report PDFs. The values are therefore entered as
  documented constants in \code{01\_collect\_data.py}, each tagged with its
  source figure number (e.g.\ AI Index 2026, Fig.~4.2.8), exactly the pattern the
  sibling project uses for the ECB survey constants. They are the project's one
  set of ``verify before publishing'' numbers; refresh them from the next report
  edition.
\end{description}

\begin{table}[htbp]\centering
\caption{Private AI investment, US vs Europe vs China (US\$bn), the anchor years
behind chart~A2. ``Europe'' is the AI Index regional line and \emph{includes the
UK}; see the comparability note below. Each row is read from the edition shown.}
\label{tab:invest}
\renewcommand{\arraystretch}{1.15}
\begin{tabular}{@{}lcccl@{}}
\toprule
year & United States & Europe & China & source edition \\
\midrule
2023 & 67.2  & 11.0 & 7.8  & AI Index 2024, Fig.~4.3.10 (line: EU+UK) \\
2024 & 109.1 & 19.4 & 9.3  & AI Index 2025, Fig.~4.3.10 (line: Europe) \\
2025 & 285.9 & 20.9 & 12.4 & AI Index 2026, Fig.~4.2.11 (line: Europe) \\
\bottomrule
\end{tabular}
\end{table}

% ----------------------------------------------------------------------------
\subsection{Two comparability caveats that matter}\label{sec:caveats}
% ----------------------------------------------------------------------------
The two halves of the thesis are measured in \emph{incommensurable units}---a
share of firms (adoption) and a flow of dollars (investment). They are
deliberately \textbf{never combined into a single ratio}; they are two parallel
lenses on the same gap, presented side by side. Beyond that, two specific caveats
shape how the investment numbers may be read:
\begin{enumerate}[label=(\roman*)]
  \item \textbf{``Europe'' includes the United Kingdom.} The AI Index regional
  time-series line bundles the UK with the EU (the 2024 edition literally labelled
  it ``European Union and United Kingdom''), and the report publishes \emph{no}
  clean EU-27 aggregate. For an EU-vs-UK split one must use the single-country
  bars (chart~A), where the UK, France, Germany, Belgium and Sweden appear
  separately.
  \item \textbf{China is understated.} The private-investment figures exclude
  state ``guidance funds''; an estimated US\$184\,bn of such funds went to Chinese
  AI firms over 2000--2023. The US/China gap on \emph{private} money therefore
  overstates the total-capital gap.
\end{enumerate}

% ============================================================================
\section{Descriptive evidence}\label{sec:descriptive}
% ============================================================================

The charts are intentionally simple; together they are the argument.

\projfig{0.86}{A_investment_by_geo}{\textbf{Leg A --- the funding gap.} Private AI
investment by geography, 2025 (Stanford AI Index; data: Quid). The United States
bar runs off the field on purpose: at US\$285.9\,bn it is $\approx 23\times$
China and $48\times$ the UK. The European leaders---the UK, France, Germany,
Belgium, Sweden---are slivers by comparison.}

\projfig{0.74}{A2_investment_gap}{\textbf{Leg A --- the gap widening.} Private AI
investment, US vs Europe vs China, 2023--2025 (Table~\ref{tab:invest}). The US
curve goes vertical while Europe and China stay pinned near the floor; since 2024
US private investment rose $\approx 160\%$, against $\approx 7\%$ for Europe.
``Europe'' includes the UK, and pre-2024 is an earlier report vintage.}

\projfig{0.66}{B_adoption_by_country}{\textbf{Leg B / C --- the adoption map.}
Share of enterprises (10+ employed) using $\ge 1$ AI technology, by EU country,
2025 (Eurostat \dataset{isoc\_eb\_ai}). A north-west/south-east split: Denmark
($42\%$) leads, Romania ($5\%$) trails---an eight-fold range. The dashed line is
the EU-27 average; colour marks an editorial Nordic-and-Western vs
Southern-and-Eastern grouping. Estonia sits in the East yet scores high.}

\projfig{0.74}{B2_adoption_trend}{\textbf{Leg B --- rising off a low base.} AI
adoption over time, EU-27 (bold) plus a country spread. The EU rate roughly
doubled in two years. The $2023\!\to\!2024$ segment is shaded because the survey
definition was broadened then; the $2024\!\to\!2025$ rise is on the consistent
definition.}

\projfig{0.60}{B3_adoption_by_firmsize}{\textbf{Leg C --- the SME gap.} EU-27 AI
adoption by firm size, 2025. Large firms ($55\%$) adopt at roughly three times the
rate of small firms ($17\%$): AI is, so far, a big-firm technology.}

% ============================================================================
\section{Interpretation}\label{sec:interpretation}
% ============================================================================

Putting the legs together:
\begin{itemize}
  \item \textbf{Leg A (investment).} Unambiguous and stark. The United States
  does not merely lead private AI investment---it laps the field
  (Fig.~\ref{fig:A_investment_by_geo}), and the lead is \emph{accelerating}
  (Fig.~\ref{fig:A2_investment_gap}): US private AI investment more than doubled
  in 2025 alone, while Europe barely moved. Cumulatively since 2013 the US has
  attracted some US\$757\,bn against Europe's low tens of billions, and in
  generative AI specifically US investment ($\approx$US\$164\,bn in 2025) dwarfs
  all of Europe's \emph{total} AI investment several times over.
  \item \textbf{Leg B (adoption).} The mirror image. European \emph{use} of AI is
  climbing fast---the EU-27 rate roughly doubled in two years to $\approx 20\%$
  (Fig.~\ref{fig:B2_adoption_trend}). Europe is a large, fast-growing customer for
  a technology it overwhelmingly does not fund or build: a consumer, not a
  producer.
  \item \textbf{Leg C (within-Europe divides).} The using is uneven on two axes.
  Geographically, a Nordic-and-Western core (Denmark $42\%$) runs far ahead of the
  Southern and Eastern periphery (Romania $5\%$)
  (Fig.~\ref{fig:B_adoption_by_country}). By size, large firms adopt at $\sim 3
  \times$ the SME rate (Fig.~\ref{fig:B3_adoption_by_firmsize}). The single EU
  ``adoption rate'' hides a wide distribution.
\end{itemize}
The synthesis is the headline: \emph{the EU is integrating AI into its economy
while the economics of \emph{owning} that AI accrue almost entirely elsewhere.}
The adoption gap is closing; the investment gap is not.

% ============================================================================
\section{Limitations and caveats}\label{sec:limits}
% ============================================================================
\begin{itemize}
  \item \textbf{Series break (adoption).} The $2023\!\to\!2024$ jump partly
  reflects a broadened survey definition, not pure growth; lead with the clean
  $2024\!\to\!2025$ step or annotate the break.
  \item \textbf{``Europe'' includes the UK (investment).} The AI Index publishes
  no clean EU-27 investment aggregate; the regional line bundles the UK
  (Sec.~\ref{sec:caveats}). EU-only figures must be built from the country bars.
  \item \textbf{China understated.} Private-investment figures omit state guidance
  funds (Sec.~\ref{sec:caveats}).
  \item \textbf{Extensive, not intensive, adoption.} \code{E\_AI\_TANY} records
  whether a firm uses \emph{any} AI, not how much value it derives; a high
  adoption rate is not the same as deep deployment.
  \item \textbf{Investment data are hand-entered.} Because no free API exposes the
  Quid series, the investment numbers are constants read from the report PDFs and
  must be re-verified and refreshed each report cycle.
  \item \textbf{Editorial geographic grouping.} The Nordic-Western vs
  Southern-Eastern colouring in leg~B is a simplification; the underlying
  continuous adoption rate is the rigorous version.
  \item \textbf{Vintages.} The adoption series is pulled live and revised by
  Eurostat; regenerate before citing any number.
\end{itemize}

% ============================================================================
\appendix
\section{Reproduction details}
% ============================================================================

\subsection{File tree}
\begin{lstlisting}
01_collect_data.py     # STEP 1  data collection -> data/*.csv
02_make_figures.py     # STEP 2  figures         -> figures/*.png
run_all.sh             # pip install + run 01 -> 02 (classic Python, no venv)
requirements.txt       # pinned dependencies
data/                  # generated tidy CSVs
figures/               # generated charts
README.md              # short project readme
documentation.tex      # this document
\end{lstlisting}

\subsection{Commands}
\begin{lstlisting}
# one-shot, fully reproducible
bash run_all.sh

# or stage by stage (after: pip install -r requirements.txt)
python 01_collect_data.py
python 02_make_figures.py
\end{lstlisting}

\subsection{Data dictionary (generated CSVs)}
\begin{center}
\renewcommand{\arraystretch}{1.2}
\begin{tabular}{@{}p{5.6cm} p{3.0cm} p{5.7cm}@{}}
\toprule
file & shape & contents \\
\midrule
\code{ai\_adoption\_by\_country.csv} & latest year & geo, value (\% firms), year \\
\code{ai\_adoption\_trend.csv} & long panel & geo, year, value (\% firms) \\
\code{ai\_adoption\_by\_firmsize.csv} & latest year & size\_emp, value (\%), year \\
\code{ai\_investment\_by\_geo.csv} & 2025 & country, investment (US\$bn), year \\
\code{ai\_investment\_us\_eu\_cn.csv} & 2023--25 & year, US, Europe, China, source \\
\code{ai\_investment\_cumulative.csv} & 2013--25 sum & country, cumulative (US\$bn) \\
\code{ai\_investment\_genai.csv} & 2025 & region, gen-AI investment (US\$bn) \\
\bottomrule
\end{tabular}
\end{center}

\vfill
\noindent\rule{\linewidth}{0.4pt}\\[2pt]
{\small The charts in this document come from a clean end-to-end run and are
reproducible via \code{run\_all.sh}. The adoption series is pulled live and
revised by Eurostat, and the investment figures are read from periodically
updated reports; regenerate the data and re-compile before citing.}

\end{document}
